\documentclass[11pt, a4paper]{article}

\usepackage[T1]{fontenc}
\usepackage[utf8]{inputenc}
\usepackage{tgpagella}
\usepackage[spanish, es-noshorthands]{babel}
\usepackage{amsmath, amssymb}
\usepackage[most]{tcolorbox}
\usepackage[margin=2.5cm]{geometry}
\usepackage{fancyhdr}

\pagestyle{fancy}
\fancyhf{}
\setlength{\headheight}{16pt}
\lhead{\textbf{UCB Sede Tarija} | Ing. Mecatrónica}
\rhead{IMT-221 | Soluciones S06-2}
\renewcommand{\headrulewidth}{0.4pt}

\begin{document}

\section*{Solucionario y Guía Docente}

\subsection*{Ejercicio 1: Diagnóstico DC}
1. $I_B = (10 - 0.7) / 470\text{k} \approx 19.79\,\mu\text{A}$. \\
2. $I_C = 100(19.79\,\mu\text{A}) \approx 1.979\,\text{mA}$. \\
3. $V_{CE} = 10 - (1.979\text{m})(2\text{k}) \approx 6.04\,\text{V}$. \\
4. $V_{BC} = 0.7 - 6.04 = -5.34\,\text{V} < 0$ (Consistente). \\
5. $Q \approx (1.98\,\text{mA}, 6.04\,\text{V})$. \\
6. Representa el estado estático de referencia en ausencia de perturbaciones externas. \\
7. Porque los parámetros incrementales ($g_m, r_\pi$) dependen directamente del valor de corriente DC en el punto $Q$.

\subsection*{Ejercicio 2: Parámetros Guiados}
1. $g_m = 2\text{mA} / 25.85\text{mV} \approx 77.37\,\text{mS}$. \\
2. $r_e \approx 1/g_m \approx 12.93\,\Omega$. \\
3. $r_\pi = 120 / 0.07737 \approx 1.55\,\text{k}\Omega$. \\
4. $g_m$ en Siemens (o mS), $r_e$ y $r_\pi$ en Ohmios. \\
5. \textbf{Predicción:} $g_m$ se duplica, $r_e$ se reduce a la mitad, $r_\pi$ se reduce a la mitad. \\
\textbf{Error común:} $r_e = R_E$. \textit{Sonda pedagógica:} "¿Cuál es el resistor físico externo?"

\subsection*{Ejercicio 4: Transformación AC}
1. $+V_{CC}$ se cortocircuita a Tierra AC. \\
2. $R_B$ permanece, ahora conectado de Base a Tierra AC. \\
3. $R_C$ permanece, conectado de Colector a Tierra AC. \\
\textbf{Error común:} Eliminar $R_B$. \textit{Sonda:} "¿Qué ocurre con la terminal que estaba atada a una fuente de tensión ideal sin variación incremental?"

\subsection*{Ejercicio 5: Ganancia Guiada}
1. $g_m = 1.5\text{mA} / 25.85\text{mV} \approx 58.03\,\text{mS}$. $r_\pi = 100 / 0.05803 \approx 1.72\,\text{k}\Omega$. \\
2. $A_v \approx -(0.05803)(2200) \approx -127.7$. \\
3. Inversión: un aumento en $v_{be}$ eleva $i_c$, aumentando la caída en $R_C$ y bajando $v_C$. \\
4. $v_o \approx (-127.7)(4\text{mV}) \approx -0.511\,\text{V}$. \\
5. La señal de gran amplitud puede acercarse a los límites de Corte o Saturación, distorsionando la onda. \\
\textbf{Error común:} Interpretar ganancia negativa como atenuación. \textit{Sonda:} "¿Qué te dice el signo respecto a la fase?"

\subsection*{Ejercicio 6: Carga Externa}
1. $g_m \approx 38.68\,\text{mS}$, $r_e \approx 25.85\,\Omega$, $r_\pi \approx 3.88\,\text{k}\Omega$. \\
2. $R_C \parallel R_L = 3.3\text{k} \parallel 10\text{k} \approx 2.48\,\text{k}\Omega$. \\
3. $A_v \approx -(0.03868)(2481) \approx -96.0$. \\
4. Suposición: Emisor atado a tierra en AC, efecto Early despreciado ($r_o \to \infty$).

\end{document}
